\documentclass[11pt]{article}
% Shared print layout. The web's provisional design and MathJax are independent.
\usepackage[utf8]{inputenc}
\usepackage[T1]{fontenc}
\usepackage[spanish,es-nodecimaldot]{babel}
\usepackage{lmodern}
\usepackage{amsmath,amssymb}
\usepackage{graphicx}
\usepackage{geometry}
\geometry{a4paper,left=20mm,right=20mm,top=25mm,bottom=22mm,headheight=16pt,headsep=8mm}
\usepackage{fancyhdr,lastpage}
\usepackage{enumitem}
\usepackage{xcolor}
\usepackage{pgfplots}
\pgfplotsset{compat=1.18}
\usepgfplotslibrary{fillbetween}
\usepackage[hidelinks,pdfencoding=auto]{hyperref}
\hypersetup{pdflang=es-ES,pdfauthor={},pdfsubject={Matemáticas Aplicadas a las Ciencias Sociales II}}
\pagestyle{fancy}
\fancyhf{}
\fancyhead[L]{\small Matemáticas Aplicadas a las Ciencias Sociales II}
\fancyhead[R]{\small 2.\textsuperscript{o} Bachillerato}
\fancyfoot[L]{\footnotesize Matrices II | Edición revisada 1.0}
\fancyfoot[R]{\footnotesize Página \thepage\ de \pageref{LastPage}}
\renewcommand{\headrulewidth}{0.4pt}
\setlength{\parindent}{0pt}
\setlength{\parskip}{0.5em}
\setlist{itemsep=0.3em,topsep=0.3em}
\newcommand{\apartado}[1]{\subsection*{#1}}
\pgfplotsset{
  revision/.style={
    width=0.94\linewidth,height=7.5cm,
    axis lines=middle,grid=major,grid style={gray!18},
    tick label style={font=\small},label style={font=\small},
    xlabel={$x$},ylabel={$y$},samples=160,
    legend style={at={(0.5,-0.16)},anchor=north,draw=none,font=\small},
    clip=true
  }
}

\hypersetup{pdftitle={Matrices II - Examen revisado}}
\begin{document}
\begin{minipage}{0.17\linewidth}
\includegraphics[width=0.80\linewidth]{logo-al-andalus.png}
\end{minipage}\hfill
\begin{minipage}{0.79\linewidth}
{\Large\bfseries Matrices II}\\[0.4em]
Duración: \textbf{55 minutos}. Puntuación: \textbf{10 puntos}.
\end{minipage}

Apellidos y nombre: \hrulefill

Fecha: \hrulefill

Responde razonadamente y numera las páginas de tu respuesta. Se permite el uso de calculadora.
\section*{Ejercicio 1 \hfill 5 puntos}
Se consideran las matrices
\[
A=\begin{pmatrix}1&1&3\\1&2&-1\\1&-1&-1\end{pmatrix},\qquad
B=\begin{pmatrix}2\\3\\0\end{pmatrix}.
\]
\begin{enumerate}[label=(\alph*)]
\item (3 puntos) Justifica que $A$ tiene inversa y calcula $A^{-1}$.
\item (2 puntos) Calcula, si existe, la matriz $X$ que satisface $AX=B$.
\end{enumerate}
\section*{Ejercicio 2 \hfill 2,5 puntos}
En este ejercicio se consideran las matrices
\[
A=\begin{pmatrix}1&-1\\0&1\end{pmatrix},\quad
B=\begin{pmatrix}2&0\\-1&2\end{pmatrix},\quad
C=\begin{pmatrix}0\\1\end{pmatrix},\quad
D=\begin{pmatrix}-1&2\end{pmatrix}.
\]
\textbf{Efectúa cuatro de las cinco operaciones} cuando sean posibles; si alguna no es posible, justifica por qué:
\[
AC+BD^t,\qquad A+BC,\qquad B^2+CD,\qquad |5A|,\qquad |C|.
\]
\section*{Ejercicio 3 \hfill 2,5 puntos}
Determina $a$ y $b$ para que se cumpla
\[
\begin{pmatrix}2&-1\\3&-1\end{pmatrix}\begin{pmatrix}a\\-b\end{pmatrix}
=\begin{pmatrix}1&2\\b&-1\end{pmatrix}\begin{pmatrix}3\\0\end{pmatrix}.
\]
\end{document}
